% ECONOMICS_PROBLEM: {"id": "Q54-163", "title": "Persistent climate damages", "jel_code": "Q54", "source_id": "OP-0430"}
% !TeX program = xelatex
\documentclass[11pt,a4paper]{article}
\usepackage[margin=23mm]{geometry}
\usepackage{fontspec}
\setmainfont{Latin Modern Roman}
\usepackage{amsmath,amssymb,mathtools}
\usepackage{xcolor,enumitem,booktabs,longtable,array,xurl,hyperref}
\definecolor{muted}{HTML}{53636C}
\hypersetup{colorlinks=true,urlcolor=blue,linkcolor=blue}
\setlength{\parindent}{0pt}
\setlength{\parskip}{5pt}
\newcommand{\R}{\mathbb R}
\newcommand{\E}{\mathbb E}
\newcommand{\N}{\mathbb N}
\newcommand{\ind}{\mathbf 1}
\newcommand{\Var}{\operatorname{Var}}
\newcommand{\Cov}{\operatorname{Cov}}
\newcommand{\supp}{\operatorname{supp}}
\DeclareMathOperator*{\argmin}{arg\,min}
\DeclareMathOperator*{\argmax}{arg\,max}
\newcommand{\entrymeta}[3]{{\sffamily\small\color{muted}\textbf{\texttt{#1}}\quad|\quad #2\par #3\par}}
\newcommand{\Problemref}[1]{\texttt{#1}}
\newcommand{\JELref}[2]{\texttt{#2} (JEL \texttt{#1})}
\begin{document}
\section*{Persistent climate damages}
\textbf{Problem ID:} \texttt{Q54-163}\quad \textbf{JEL:} \texttt{Q54}\par
% BEGIN ECONOMICS STATEMENT
\label{problem:OP-0430}
{\sffamily\small\color{muted}Original subject group: Climate, environment and energy\par}
\label{definitions:problem:30007712}
\textbf{Audit classification:} Background; proposed extension. Required a complete climate intervention and formalized finite persistence tests.
\subsection*{Definitions and precise research target}
\textit{Editorial formalization.} Consider the countries with annual national accounts and temperature records from 1960 through 2025. For country \(i\), let \(Y_{it}(\delta)\) be real output per resident in year \(t\) under a climate intervention that permanently raises its expected annual temperature by \(\delta\) degrees Celsius from a common initial year. The intervention preserves the baseline distribution of nonclimate shocks and permits households, firms, and governments to adapt; migration and capital accumulation are therefore outcomes. Let \(Y_{it}(0)\) denote the corresponding baseline. For every horizon \(h\in\{1,\ldots,50\}\), define
\[\theta_h=\left.\frac{\partial}{\partial\delta}E[\log Y_{i,t+h}(\delta)-\log Y_{i,t+h}(0)]\right|_{\delta=0}.\]
Assume output is positive and the displayed mean derivative exists. The expectation uses fixed baseline population weights across countries. The question is whether a credible identification strategy and economically disciplined extrapolation can distinguish an approximately constant response \(\theta_h\) after a prespecified lag from a response with persistently negative first differences over the stated horizon. Predeclare a tolerance for approximate constancy; every finite sequence is bounded, so boundedness cannot distinguish these cases. Weather fluctuations alone need not identify this intervention. Specify restrictions connecting temporary temperature shocks to permanent climate changes, account for cross-country spillovers, and report uncertainty when those restrictions are relaxed. The target is this finite-horizon response, not a universal claim that every warming episode permanently reduces growth.

Fix a baseline date \(t_0\), country set \(I\), and weights \(\omega_i\geq0\), \(\sum_i\omega_i=1\). The climate intervention must specify the joint temperature, precipitation, and extremes process, not only its mean: for example change one declared global physical parameter and map it to every country. Define \(g_h(\delta)=\sum_i\omega_iE[\log Y_{i,t_0+h}(\delta)]\) and \(\theta_h=g_h'(0)\). Assume positive output and integrability on a neighborhood of zero; differentiability of the expectation is an assumption, not an automatic exchange of derivative and expectation. Fix lag \(h_0<50\) and \(\epsilon>0\): level persistence means \(\max_{h\geq h_0}|\theta_h-\theta_{h_0}|\leq\epsilon\); continuing growth loss means \(\theta_{h+1}-\theta_h\leq-2\epsilon\) for all \(h=h_0,\ldots,49\). These classes need not exhaust possible response profiles. All expectations use a stated baseline population and probability law. Identification means every admissible data-generating model with the same observed-data law yields the same displayed target; otherwise report the identified set. Exact equations, horizons and assumptions are editorial, rather than source-given canonical conjectures.
\subsection*{Variants and assumption changes}
\begin{enumerate}
\item \textbf{Editorial temporary-shock variant.} Replace permanent warming by a one-year temperature perturbation at \(t_0\); define \(\theta_h^{\rm pulse}=\partial_\delta E[\log Y_{t_0+h}(\delta)-\log Y_{t_0+h}(0)]|_0\). Compare permanent and pulse responses only under a declared linear time-invariant response law.
\item \textbf{Editorial adaptation variant.} Introduce constrained and unconstrained adaptation regimes \(a\in\{0,1\}\), with common climate paths. Identify \(\theta_h^1-\theta_h^0\); holding migration/capital fixed changes the target.
\end{enumerate}
\subsection*{References and source audit}
\textbf{1.} Climate macroeconomics survey; no exact fifty-year or tipping-event conjecture verified. Locator: NBER working-paper record and abstract; AEA forthcoming record, DOI 10.1257/jel.20261831. Primary indexed record verified; direct NBER/PDF retrieval failed. AEA forthcoming abstract inspected; no final journal publication year asserted. URL: \url{https://www.nber.org/papers/w33567}.

\textbf{2.} SCC synthesis discusses damage persistence, tipping, and nonmarket substitution; proposed estimands are editorial. Locator: Introduction; Section 1.1; Section 2. Full text or primary publication page inspected. URL: \url{https://pmc.ncbi.nlm.nih.gov/articles/PMC11670083/}.
\begin{enumerate}\raggedright
\item\hypertarget{definitions:source:30007712:1}{} Adrien Bilal; James H. Stock. 2025. \emph{A Guide to Macroeconomics and Climate Change}. NBER working-paper record and abstract; AEA forthcoming record, DOI 10.1257/jel.20261831.
\url{https://www.nber.org/papers/w33567}.
DOI: \url{https://doi.org/10.3386/w33567}.
\item\hypertarget{definitions:source:30007712:2}{} Frances C. Moore; Moritz A. Drupp; James Rising; Simon Dietz; Ivan Rudik; Gernot Wagner. 2024. \emph{Synthesis of evidence yields high social cost of carbon due to structural model variation and uncertainties}. Introduction; Section 1.1; Section 2.
\url{https://pmc.ncbi.nlm.nih.gov/articles/PMC11670083/}.
DOI: \url{https://doi.org/10.1073/pnas.2410733121}.
\end{enumerate}

\subsection*{Common conventions: Source mathematical conventions}

These conventions apply to each entry unless explicitly replaced.

\textbf{Models and identification.} A primitive model \(m\) specifies all probability laws, feasible choices, information, equilibrium and measurement rules used by its target. The admissible class \(\mathcal M\) is fixed independently of outcomes. If \(P_O(m)\) is its observable probability law and \(\tau(m)\) the target, its identified set at observable law \(P\) is
\[
 \mathcal I_\tau(P)=\{\tau(m):m\in\mathcal M,\ P_O(m)=P\}.
\]
Point identification means that this set is a singleton. An empty set signals incompatibility of data and assumptions. Coordinate bounds need not describe the joint set. Bounds are sharp only if every retained target is attainable or arbitrarily approachable by compatible models. Finite bounds require explicit restrictions on unobserved quantities; unrestricted confounding, hidden wealth, extrapolation or arbitrary equilibrium selection can yield uninformative sets.

\textbf{Causal interventions.} A potential outcome \(Y_i(p)\) is the outcome under a complete policy regime \(p\), including timing, financing, information and market spillovers. The target population and units are fixed before treatment. With interference, use the complete assignment vector or a justified exposure mapping; individual no-interference notation alone cannot identify an economy-wide intervention. A difference of mean potential outcomes does not require the cross-world joint distribution. Conditioning on post-treatment migration, survival, enrollment or employment changes the estimand and needs separate justification.

\textbf{Statistical statements.} An estimator, confidence set, forecast or test specifies a sampling law, model class and loss/coverage criterion. Uniform \((1-\alpha)\)-coverage means \(\inf_{m\in\mathcal M}\Pr_m\{\tau(m)\in C_n\}\geq1-\alpha\), or an explicitly asymptotic version. The whole parameter space trivially has coverage, so informativeness requires a stated width, risk or power objective. A forecasting improvement is relative to a named benchmark, information set and evaluation population.

\textbf{Equilibrium and optimization.} An equilibrium comprises feasible plans, optimality, beliefs where needed, budget constraints and all market/resource clearing. The primitives or a stated benchmark define each correspondence \(\mathcal E(p;m)\). Nonemptiness is assumed or proved, never inferred just from compact policy sets. A selected-equilibrium target uses a declared selection rule; a robust target quantifies over all permitted equilibria. Use suprema or infima unless attainment is established. An \(\varepsilon\)-optimal policy is feasible and within \(\varepsilon\) of the finite optimum value. Report an empty feasible set separately.

\textbf{Welfare and accounting.} Normative utility, interpersonal normalization, social weights, numeraire, discounting and the use of public receipts are primitives. Behavioral utility need not coincide with the welfare criterion. Household welfare includes ownership income and rents once; adding profits again to compensating variation generally double-counts them. A monetary equivalent is defined by an explicit transfer and price convention, with monotonicity and range conditions. Average effects, mechanism contributions, transfers and welfare losses are different targets. Infinite sums require convergence and dynamic feasible plans require terminal or no-Ponzi conditions.

\textbf{Source versions.} A working-paper date, revision date and journal publication year are distinguished. A DOI for a journal version is not automatically the DOI of an earlier manuscript. Locators refer to the stated version and printed pages where specified. A metadata-only check does not verify a claim about a full-text conclusion. Every entry's source subsection narrows the provenance claim accordingly.
% END ECONOMICS STATEMENT
\end{document}
